\chapter{断层预测}

\section{任务背景}

断层预测的目标是从三维地震振幅体中识别不连续构造。输入为局部地震体
$\bm{X}\in\mathbb{R}^{H\times W\times D}$，输出为逐体素断层概率
$\bm{P}\in[0,1]^{H\times W\times D}$。断层通常只占很少的体素，类别极不平衡；
同时，断层面具有跨切片延续性，单张切片上的局部纹理不足以完整描述空间结构。

本赛道先建立可解释的局部逻辑回归基线，再评估 SAM-Med3D 的三维预训练表征
\cite{wang2023sammed3d}。研究问题不是简单判断“大模型是否更强”，而是检验三维
预训练表征能否在合法开发折上稳定改善断层分割。

\section{方法原理}

基线模型从地震振幅、局部梯度和不连续性响应中构造特征向量 $\bm{z}_i$，并用
逻辑函数估计第 $i$ 个体素属于断层的概率：
\begin{equation}
  p_i^{\mathrm{base}}=\sigma(\bm{w}^{\mathsf T}\bm{z}_i+b),
  \qquad
  \sigma(a)=\frac{1}{1+\exp(-a)} .
\end{equation}
其中 $\bm{w}$ 和 $b$ 分别为特征权重与截距。该模型参数少，便于检查局部特征与
断层概率之间的关系，适合作为小样本条件下的基线。

候选分支把三维地震块送入 SAM-Med3D 编码器，得到预训练特征
$\bm{h}_i^{\mathrm{sam}}$。适配头将其映射为候选概率 $p_i^{\mathrm{sam}}$，再由
门控系数 $g_i\in[0,1]$ 与基线融合：
\begin{equation}
  p_i=(1-g_i)p_i^{\mathrm{base}}+g_i p_i^{\mathrm{sam}} .
\end{equation}
当开发集证据不足时令 $g_i=0$，模型自动退回基线。这一设计把候选表征与默认预测
分开，避免未经验证的三维分支直接改变结果。

\ArchitectureFigure{fault}
  {断层预测模型架构。局部可解释基线与 SAM-Med3D 三维表征通过受控门融合。}
  {fault-architecture}

\Cref{fig:fault-architecture}描述目标模型的数据流。方法图中的三维分支能否进入
定量比较，仍取决于是否存在连续三维开发体。

\section{训练策略}

训练采用类别加权二元交叉熵。设正样本权重为 $\omega_1$，负样本权重为
$\omega_0$，则损失为
\begin{equation}
  \mathcal{L}_{\mathrm{WBCE}}
  =-\frac{1}{N}\sum_{i=1}^{N}
  \left[\omega_1 y_i\log p_i+\omega_0(1-y_i)\log(1-p_i)\right].
\end{equation}
权重由训练集类别频率确定，本次实验的负类与正类权重分别为 $0.5053$ 和
$47.7847$。模型训练 120 个 epoch，批量大小为 16，并根据验证集损失选择第 16
个 epoch 的参数。

概率阈值也只在验证集上选择。最终阈值为 $0.4506$，对应验证集 F1 的最大值。
测试集只在模型和阈值固定后评价一次。

\TrainingFlow
  {地震体与断层标签}
  {局部块采样}
  {加权损失训练}
  {验证集选模型}
  {固定阈值测试}
  {断层预测训练流程}
  {fault-training}

\section{实验设计}

基线实验使用 198 个训练样本、56 个验证样本和 96 个测试样本。现有样本是彼此独立
的二维剖面，输入通道形状为 $1\times33\times65$，输出为逐像素断层概率。理论任务
目标仍是三维体素预测；候选实验需要在连续三维数据上增加预训练编码器，并由独立
开发折评价。

进一步检查发现，当前三维赛道没有满足约束的合法开发折。因此，SAM-Med3D 分支
尚不能完成公平排名，冻结测试集也没有用于候选模型选择。本章报告的定量结果来自
已完成独立复核的基线；方法架构用于说明后续实验的比较对象。

\section{评估指标}

断层预测同时评价检出能力和空间重叠。精确率与召回率分别为
\begin{equation}
  \Metric{Precision}=\frac{TP}{TP+FP},
  \qquad
  \Metric{Recall}=\frac{TP}{TP+FN}.
\end{equation}
Dice 系数和交并比定义为
\begin{equation}
  \Metric{Dice}=\frac{2TP}{2TP+FP+FN},
  \qquad
  \Metric{IoU}=\frac{TP}{TP+FP+FN}.
\end{equation}
由于正类稀少，准确率会被大量背景体素主导，因此本赛道以 Dice、IoU、平均精确率
（AP）和 PR-AUC 为主要指标，并用召回率检查漏检风险。

\section{实验结果}

\begin{table}[htbp]
  \centering
  \caption{断层预测基线的测试结果}
  \label{tab:fault-results}
  \begin{tabular}{lrrrrr}
    \toprule
    模型 & Precision & Recall & Dice/F1 & IoU & AP \\
    \midrule
    局部逻辑回归 & 0.0078 & 0.9022 & 0.0155 & 0.0078 & 0.0069 \\
    \bottomrule
  \end{tabular}
\end{table}

\Cref{tab:fault-results}显示，基线取得 $0.9022$ 的召回率，但精确率只有
$0.0078$。这说明模型能够覆盖大部分
真实断层，却产生了大量误检。Dice 和 IoU 均较低，当前模型更适合作为高召回的
候选区域生成器，尚不足以直接承担精确断层制图。

\ResultFigure{0.92}{fault}{research_geological_context.png}
  {Volve 断层棒与 BCU 层位在真实 UTM--TWT 坐标中的空间关系。该图展示解释几何，
  不把稀疏断层棒扩展为未验证的连续概率体。}

\ResultFigure{0.94}{fault}{research_seismic_overlay.png}
  {Inline 10243 的原始地震剖面与同一剖面上的断层解释控制点。连线仅连接同一断层棒
  中至少两个控制点，用于核对解释位置与反射不连续结构。}

三维场景说明断层解释并不是无坐标的样本集合，而是位于已配准地震网格中的空间
结构。剖面核对进一步显示，进入连线的控制点主要分布在反射同相轴发生弯折或中断的
区域。两幅图只验证解释与地震背景的几何一致性，不能替代模型的体素级预测评价。

\ResultPair{fault}
  {real_test_panel.png}{测试样本上的输入、标签与预测}
  {spatial_context.png}{断层预测的空间位置与邻域结构}
  {断层预测的样本级与空间级结果}

\ResultPair{fault}
  {loss_curve.png}{训练与验证损失}
  {prediction_visualization.png}{概率预测与二值结果}
  {基线模型的优化过程与预测表现}

当前实验的主要问题是误检，而不是漏检。后续实验应先补齐合法三维
开发折，再比较局部基线、随机初始化三维编码器和 SAM-Med3D 三组模型；只有当预训练
分支在独立开发折上稳定提高 Dice、PR-AUC 和表面距离指标时，才有理由进入最终预测。

\FloatBarrier
